\section{Reproducible implementation and discriminating experiments}
\label{sec:C03:experiments}
An implementation record should separate architecture configuration from model weights, tokenizer artifacts, and policy-wrapper code. The configuration specifies head counts, sharing patterns, dimensions, positional treatment, context limits, and expert routing where applicable. The tokenizer record specifies vocabulary, merge rules or learned model, normalization, special tokens, and decoder behavior. The wrapper record specifies message serialization, grammar masks, parser handling, retries, and routing. A single model name rarely fixes all three. Reproduction requires immutable artifacts or hashes and a statement of any unavailable component.

The appropriate experiment is a factorial comparison rather than a universal leaderboard. Hold a checkpoint fixed while varying tokenizer-compatible serialization, schema ordering, constrained versus unconstrained decoding, and parser behavior. Separately compare model checkpoints under a common wrapper when their interfaces permit it. Include exact-copy tasks for identifiers, multilingual instructions, code editing, and multi-step tools with changing state. The outcome set should include useful task completion, rejected operations, invalid arguments, permission violations, and resource consumption. These measurements identify whether a change improves language modeling, interface compatibility, or the execution policy.

For architecture comparisons, define the resource axis before selecting models. Matching total parameters favors some sparse architectures; matching active FLOPs ignores resident memory; matching wall-clock time depends on implementation maturity. Report several axes when feasible and avoid declaring a single comparison intrinsically fair. A recurrent model with external retrieval should be compared with attention models using the same retrieval service, plus an ablation without retrieval. Otherwise memory infrastructure can conceal the architecture's actual failure mode. Token budgets should be translated into underlying bytes and evidence units as well as reported in native tokens.

Contamination controls extend to tools and interfaces. A benchmark can leak its expected API names or canonical action traces even if final answers are held out. Split by task template, environment instance, and repository lineage where those structures matter. Randomize harmless tool ordering to test sensitivity while retaining a fixed canonical order for the main result. Use paired task instances and repeated runs to distinguish stable gains from sampling variation. A changed model response format should not silently invalidate the baseline parser; both systems need a documented, equally capable integration.

The expected failure taxonomy is mechanistic. Segmentation failures corrupt exact strings; serialization failures cross message boundaries; contextual failures omit relevant constraints; routing failures select a wrong operation; authorization failures execute an impermissible one; semantic failures produce a valid but ineffective action. These categories can co-occur and should not be forced into a single label when evidence is ambiguous. A retained raw request and observed transition allow later reclassification. Without that evidence, a reported model mistake may actually be an execution race or a lossy context transformation.

